\chapter*{Abstract}
\addcontentsline{toc}{chapter}{Abstract}

Embedded artificial intelligence for autonomous robots lies at the intersection of sensing physics, signal processing, estimation, machine learning, computer architecture, real-time systems, and control. A paper that focuses on an event camera, a LiDAR--camera fusion network, an FPGA runtime controller, or a tightly coupled odometry estimator may appear isolated when read only at the algorithm level. These works become connected when they are placed within a common systems model that follows information from the physical world through sensing, representation, alignment, fusion, perception, state estimation, prediction, planning, control, and action.

This handbook develops that systems model and defines the foundational terminology required to use it. It explains the distinction between sensors and modalities; the properties of RGB cameras, event cameras, LiDAR, radar, inertial sensors, and global-navigation systems; the role of images, point clouds, event graphs, feature maps, occupancy maps, and state vectors; and the need for temporal synchronization, spatial calibration, coordinate transformations, and data association. Sensor fusion is organized into data-level, feature-level, decision-level, state-estimation, and runtime-management forms. Loose and tight coupling are compared through measurement models and residual-based state estimation. CNNs, GNNs, asynchronous GNNs, transformers, attention mechanisms, and spiking neural networks are positioned according to the data structures and deployment constraints they address.

The document also emphasizes the engineering dimensions that distinguish embedded AI from unconstrained cloud inference: end-to-end latency, deterministic deadlines, power, energy per decision, memory movement, communication cost, thermal limits, fault tolerance, uncertainty, and graceful degradation. Hardware platforms and co-design techniques are reviewed, including microcontrollers, GPUs, FPGAs, NPUs, compute-in-memory, neuromorphic systems, quantization, sparsity, streaming dataflow, near-sensor filtering, and runtime reconfiguration.

Finally, six Week 7 papers are analyzed as related case studies: low-latency RGB--event detection, context-aware FPGA sensor fusion, RGB--event compute-in-memory, fusion-aware near-sensor filtering, LiDAR--inertial--event--visual odometry, and attention-based LiDAR--camera fusion. Their research questions and technical relationships are mapped onto the sensing-to-decision chain, providing a structured foundation for deeper paper review and future embedded-AI research.
